\medskip

\textbf{PARTIE A}

\medskip

On considère l'équation différentielle 
$(E) : \quad y' +\dfrac14y = 20\e^{-\frac14 x},$

d'inconnue $y$, fonction définie et dérivable sur l'intervalle $[0~;~ +\infty[$.

\medskip

\begin{enumerate}
\item 
Soit $g$ la fonction  définie sur l'intervalle $[0~;~+\infty[$ par $g(x) = ax\e^{-\frac14 x}$.

La fonction linéaire $u$ définie par $u(x) = - \dfrac14 x$ est dérivable sur $\R$ et sur cet intervalle $u'(x) = - \dfrac14$.

La fonction composée $g(u(x)) = ax\e^{u(x)}$ est elle aussi dérivable sur $\R$ et sur cet intervalle, $g'(u) = a\e^{u(x)} + axu' \times \e^{u(x)} = a\e^{-\frac14 x} + ax \times \left(- \dfrac14 \right)\e^{-\frac14 x} = \e^{-\frac14 x}\left(a - \dfrac{ax}{4}\right)$.

Or $g'(x) + \dfrac14g(x) = \e^{-\frac14 x}\left(a - \dfrac{ax}{4}\right) + \dfrac14 \times ax\e^{-\frac14 x} = \e^{-\frac14 x}\left(a - \dfrac{ax}{4} + \dfrac{ax}{4} \right) = a\e^{-\frac14 x}$.

Donc $g$ est solution de l'équation différentielle sur $[0~;~ +\infty[$, si sur cet intervalle :

$a\e^{-\frac14 x} = 20\e^{-\frac14 x} \iff a = 20$, car quel que soit $x$ réel, $\e^{-\frac14 x} \ne 0$.

La fonction $g : x \longmapsto 20x \e^{-\frac14 x}$ définie sur l'intervalle $[0~;~ +\infty[$ est une solution particulière de l'équation $(E)$.

\item On considère l'équation différentielle 
$(E') : \quad y' +\dfrac14 y = 0,$

d'inconnue $y$, fonction définie et dérivable sur l'intervalle $[0~;~ +\infty[$.


L'équation $(E')$ peut s'écrire $y' = - \dfrac14y$ et on sait que les solutions de cette équation $(E')$ sont les fonctions $f$ définies par:

$f(x) = K\e^{-\frac14 x}$, avec $K \in \R$.

\item 
Les solutions $f$ de l'équation $(E)$ sont telles que pour $x \in [0~;~+ \infty[$, 

$f'(x) + \dfrac14f(x) = 20\e^{- \frac14 x}$ ou d'après la question 1. :

$f'(x) + \dfrac14f(x) = g'(x) + \dfrac14g(x) \iff f'(x) - g'(x) = \dfrac14[f(x) - g(x)]$ soit par linéarité de la dérivation :

$(f - g)'(x) + \dfrac14[(f - g ](x) = 0$ : ceci signifie que la fonction $f - g$ est solution de $(E')$ c'est-à)dire que 

$f(x) - g(x) = K\e^{-\frac14 x} \iff f(x) = K\e^{-\frac14 x} + g(x)$ ou enfin 
$f(x) = K\e^{-\frac14 x} + 20x\e^{-\frac14 x}\iff f(x) = (20x + K)\e^{-\frac14 x}$, avec $K \in \R$.

Les solutions de $(E)$ sur l'intervalle $[0~;~ +\infty[$ sont donc les fonctions $f$ définies par:\\
$f(x) = \e^{-\frac14 x}(20x + K)$, avec  $K \in \R$.

\item 
On a $f(0) = 8 \iff (20 \times 0 + K)\e^{0} = 8 \iff K = 8$, donc :

$f(x) = (20x + 8)\e^{-\frac14 x} = 4\left (5x + 2\right )\e^{-\frac14 x}$.
\end{enumerate}

\bigskip

\textbf{PARTIE B}

\medskip

On considère la fonction $f$ définie sur $[0~;~ +\infty[$ par 
$f(x) = (20x + 8)\e^{-\frac14 x}$.





\begin{enumerate}
\item 
	\begin{enumerate}
		\item 


La fonction $f$ est la fonction trouvée à la fin de la partie A : elle est donc dérivable et vérifie l'équation $(E)$, soit 

$f'(x) + \dfrac14 f(x) = 20\e^{-\frac14 x}$ donc \\
$f'(x) = 20\e^{-\frac14 x} - \dfrac14 \times 4(5x + 2)\e^{-\frac14 x} = \e^{-\frac14 x}(20 - 5x - 2) = (18 - 5x)\e^{-\frac14 x}$.

		\item 
		Comme quel que soit $x$, \, $\e^{-\frac14 x} > 0$, le signe de $f'(x)$ est celui de $18 - 5x$ :
		
$\bullet~$ $18 - 5x > 0 \iff \dfrac{18}{5} > x$ : sur l'intervalle $\left[0~;~\dfrac{18}{5}\right], \, f'(x) > 0$ : la fonction $f$ est croissante ;

$\bullet~$ $18 - 5x < 0 \iff \dfrac{18}{5} < x$ : sur l'intervalle $\left[\dfrac{18}{5}~;~+ \infty\right], \, f'(x) < 0$ : la fonction $f$ est décroissante ;

$f'\left(\dfrac{18}{5}\right) = 0$, donc $f\left(\dfrac{18}{5}\right)$ est le maximum de la fonction $f$ sur l'intervalle $[0~;~ +\infty[$.

On a $f\left(\dfrac{18}{5}\right) = (72 + 8)\e^{-\frac14 \times \frac{18}{5}} = 80\e^{-\frac{9}{10}} \approx 32,53$.

On a vu à la partie A que $f(0) = 8$, et on admet que $\displaystyle\lim_{x \to + \infty}f(x) = 0$. 

D'où le tableau de variations :

\begin{center}
{\renewcommand{\arraystretch}{1.3}
\psset{nodesep=3pt,arrowsize=2pt 3}
\def\esp{\hspace*{2.5cm}}
\def\hauteur{20pt}
$\begin{array}{|c|*5{c}|}
\hline
x & 0  & \esp & \frac{18}{5} & \esp & +\infty \\ 
\hline
f'(x) &  &   \pmb{+} & \vline\hspace{-2.7pt}0 & \pmb{-} & \\ 
\hline
 & &  &   \Rnode{max}{80 \e^{-\frac{9}{10}}}  &  &   \\  
f & &     &  &  &  \rule{0pt}{\hauteur} \\ 
 & \Rnode{min1}{8} &   &  &  &   \Rnode{min2}{0} \rule{0pt}{\hauteur}    
 \ncline{->}{min1}{max} 
 \ncline{->}{max}{min2} 
 \\ 
\hline
\end{array} $
}
\end{center}	

	\end{enumerate}
\item Dans cette question on s'intéresse à l'équation $f(x) = 8$.
	\begin{enumerate}
		\item 
		On a $f(14) = 288\e^{-3,5} \approx 8,7$ et $f(15) = 300\e^{-3.75} \approx 7,2$.
		
Sur l'intervalle [14~;~15], la fonction $f$ est continue car dérivable sur $[0~;~+ \infty[$, donc sur [14~;~15] et elle est strictement décroissante; comme $8 \in [f(15)~;~f(14)]$, il existe d'après le corollaire du théorème des valeurs intermédiaires un réel unique $\alpha$ telle que $f(\alpha) = 8$.
		\item On complète le tableau ci-dessous en faisant tourner étape par étape la fonction \texttt{solution\_equation} écrite en langage Python


\begin{center}
\begin{tabularx}{0.8\linewidth}{|c|c|*{4}{>{\centering \arraybackslash}X|}}
\hline
	$a$		&14		&14		&14,25	&14,375	&14,4375\\ \hline
	$b$		&15		&14,5	&14,5	&14,5	&14,5\\ \hline
	$b - a$	&1		&0,5	&0,25	&0,125	&0,0625\\ \hline
	$m$		&14,5	&14,25	&14,375	&14,4375	&\cellcolor[gray]{0.5}\\ \hline
\parbox{2cm}{\centering{}Condition\\ $f(m) > 8$} \rule[-12pt]{0pt}{30pt} &FAUX	&VRAI	&VRAI	&VRAI	&\cellcolor[gray]{0.5}\\ \hline
\end{tabularx}
\end{center}





















		\item 
L'objectif de la fonction \texttt{solution\_equation} est de déterminer, par dichotomie, un encadrement d'amplitude $0,1$ de la solution de l'équation $f(x)=8$ dans l'intervalle [14\, ;\,15].
	\end{enumerate}
\end{enumerate}

\vspace{0,5cm}
